% ch-a5-04.tex —— 第 5 篇（Ben-Zvi–Chen–Helm–Nadler）末片：
% 附录 A 余部（A.1.2–A.2.5）+ 参考文献 + 四位作者通讯地址。
% 底本 PDF 470–483 = 印刷页 457–470。本篇（第 5 篇）止于印刷页 470。
% 本片自印刷页 457 页首（A.1.2 小标题）起；不设章标题；参考文献条目保留英文。
\ysrc{457}

\subsubsection*{A.1.2.\quad 圆作用与 $B\GG_a$-作用的重整化}

我们接着讨论对不变量操作作重整化的必要性（这一点在 [Pr15] 中已经提出）。本小节的讨论可概括如下。
\begin{enumerate}[(1)]
\item 一个 $\mathrm{QC}(S^1)$-模范畴 $\mathbf C$ 等变化为一个循环形变 $\mathbf C^{S^1}$，它位于 2-移位形式仿射线 $\widehat{\mathbb A}^1[2]$ 之上，其特殊纤维是由 $S^1$-不变对象生成的全子范畴。
\item 一个 $\mathrm{QC}(B\GG_a)$-模范畴 $\mathbf C$ 等变化为一个循环形变 $\mathbf C^{B\GG_a}$，它位于 2-移位形式仿射线 $\widehat{\mathbb A}^1[2]$ 之上，其特殊纤维还原出 $\mathbf C$。
\item 一个 $\mathrm{QC}(B\GG_a\rtimes\GG_m)$-模范畴 $\mathbf C$ 等变化为一个循环形变 $\mathbf C^{B\GG_a\rtimes\GG_m}$，它位于分次形式仿射线 $\widehat{\mathbb A}^1/\GG_m$ 之上，其特殊纤维还原出 $\mathbf C$。
\end{enumerate}

上述各情形中的一般纤维都消失。重整化的目标，是把形式仿射线换回通常的仿射线，从而得到有意义的一般纤维。

我们先讨论一些一般性事实。设 $G$ 是作用于预叠 $X$ 的群预叠。于是 $\Perf(X)$ 上有自然的 $\Perf(G)$-余作用，$\mathrm{QC}(X)$ 上有自然的 $\mathrm{QC}(G)$-余作用。此外还有
\[
\Perf(X)^{\Perf(G)}\simeq\Perf(X/G)\subset\mathrm{QC}(X/G)\simeq\mathrm{QC}(X)^{\mathrm{QC}(G)}.
\]

若 $X/G$ 是 [BFN10] 意义下的完美叠，即 $\mathrm{QC}(X/G)$ 由 $\Perf(X/G)$ 紧生成，那么先取小范畴的 $G$-不变量再作 ind-完备化，与取大范畴的 $G$-不变量相同。据 [BFN10] 的推论 3.22，例如在特征零、$X$ 拟射影而 $G$ 是仿射代数群时，情况就是这样。然而，我们的群 $G=S^1,B\GG_a$ 越出了这一框架：如例 A.5 所见，$BG$ 并不完美。由于叠 $BS^1$ 与 $B^2\GG_a$ 都不完美，下列两种等变化会给出不同的答案：
\begin{enumerate}[(1)]
\item 取大范畴的 $\mathrm{QC}(G)$-不变量，正如上一节所做的那样；
\item 取小范畴的 $\Perf(G)$-不变量，再作 ind-完备化，得到一个大范畴。
\end{enumerate}

我们现在采用第二条路线，并定义如下重整化等变化。

\begin{defn}{A.7}\label{defn:a5-04-1}
设 $G$ 是预叠中的群对象。我们用 $\mathbf C^G$ 记余模范畴的 $\mathrm{QC}(G)$-不变量。设 $\mathbf C=\Ind(\mathbf C_0)$ 是紧生成范畴，且 $\mathbf C_0$ 带 $\Perf(G)$-余作用。我们定义\emph{紧重整化 $G$-不变量范畴}为 $\Ind\Perf(BG)$-模范畴
\[
\mathbf C^{\omega G}:=\Ind\bigl(\mathbf C_0^{\Perf(G)}\bigr).
\]
\end{defn}

由函子 $\mathbf C_0^{\Perf(G)}\to\mathbf C^{\mathrm{QC}(G)}$ 诱导出自然函子 $\mathbf C^{\omega G}\to\mathbf C^{G}$。

\ysrc{458}

于是，对我们所关心的几种情形 $G=S^1,B\GG_a,B\GG_a\rtimes\GG_m$，紧重整化等变化具有如下形式。
\begin{enumerate}[(1)]
\item 一个 $\Ind\Perf(S^1)$-模范畴 $\mathbf C$ 紧等变化为 2-移位仿射线 $\mathbb A^1[2]$ 上的循环形变 $\mathbf C^{\omega S^1}$。
\item 一个 $\Ind\Perf(B\GG_a)=\mathrm{QC}(B\GG_a)$-模范畴 $\mathbf C$ 紧等变化为 2-移位仿射线 $\mathbb A^1[2]$ 上的循环形变 $\mathbf C^{\omega B\GG_a}$。
\item 一个 $\Ind\Perf(B\GG_a\rtimes\GG_m)=\mathrm{QC}(B\GG_a\rtimes\GG_m)$-模范畴 $\mathbf C$ 紧等变化为分次仿射线 $\mathbb A^1/\GG_m$ 上的循环形变 $\mathbf C^{\omega B\GG_a\rtimes\GG_m}$。
\end{enumerate}

由这一计算，我们可以作如下定义。

\begin{defn}{A.8}\label{defn:a5-04-2}
设 $\mathbf C=\Ind(\mathbf C_0)$ 是紧生成范畴，带有 $\mathrm{QC}(S^1)$-作用，该作用限制在紧对象上。我们定义 \emph{Tate 构造}为 2-周期范畴
\[
\mathbf C^{\mathrm{Tate}}:=\mathbf C^{\omega S^1}\otimes_{\Mod(k[u])}\Mod(k[u,u^{-1}]).
\]
对带有 $\mathrm{QC}(B\GG_a)$-作用的范畴可作类似定义。若范畴带 $\mathrm{QC}(B\GG_a\rtimes\GG_m)$-作用，则可定义\emph{分次 Tate 构造}
\[
\mathbf C^{\mathrm{Tate}\rtimes\GG_m}:=\mathbf C^{\omega(B\GG_a\rtimes\GG_m)}\otimes_{\mathrm{QC}(\mathbb A^1/\GG_m)}\mathrm{QC}(\GG_m/\GG_m),
\]
它是（非周期的）$k$-线性范畴。
\end{defn}

在 [BCHN] 中我们证明：紧重整化不变量范畴与通常不变量在 $u=0$ 纤维处相容。

\begin{prop}{A.9}\label{prop:a5-04-3}
设 $\mathbf C_0$ 是小稳定 $k$-线性 $\infty$-范畴，带 $\Perf(S^1)$-余作用，并令 $\mathbf C=\Ind(\mathbf C_0)$。则典范函子
\[
\mathbf C^{\omega S^1}\otimes_{k[u]}\mathrm{Vect}_k\isom\mathbf C^{S^1}\otimes_{k[u]}\mathrm{Vect}_k
\]
是等价。特别地，由命题 A.6，函子
\[
\mathbf C^{\omega S^1}\otimes_{k[u]}\mathrm{Vect}\into\mathbf C
\]
全忠实。对紧重整化的 $B\GG_a$-不变量与 $B\GG_a\rtimes\GG_m$-不变量，类似的陈述也成立。
\end{prop}

\subsubsection*{A.1.3.\quad Koszul 对偶与分次提升}

正如定义 A.8 所示，Tate 构造给出 2-周期范畴，而分次 Tate 构造给出到非周期范畴的提升。这些现象是 Koszul 对偶 [BGS96,BY13] 的典型特征；下面我们给出这些概念的形式化。

\begin{defn}{A.10}\label{defn:a5-04-4}
设 $\mathbf C$ 是 $k$-线性范畴。$\mathbf C$ 的一个\emph{分次提升}，指一个 $\mathrm{QC}(B\GG_m)$-模范畴 $\mathbf C_{gr}$ 以及一个等价 $\mathbf C\simeq\mathbf C_{gr}\otimes_{\mathrm{QC}(B\GG_m)}\mathrm{Vect}_k$。$\Z$-分次向量空间范畴 $\mathrm{QC}(B\GG_m)$ 有一个由 Tate 剪切给出的自同构：
\[
M^{\sslash}:=\bigoplus_{n\in\Z}M_n[-2n],\qquad
M^{\backslash\backslash}:=\bigoplus_{n\in\Z}M_n[2n].
\]
对任意 $\mathrm{QC}(B\GG_m)$-模范畴 $\mathbf C$，我们用 $\mathbf C^{\sslash}$（相应地 $\mathbf C^{\backslash\backslash}$）记 $\mathrm{QC}(B\GG_m)$ 此时经由 Tate 剪切（相应地反剪切）作用的那个 $\mathrm{QC}(B\GG_m)$-模范畴。对任意 $\mathrm{QC}(\GG_m)$-模范畴 $\mathbf C$，我们用 $\mathbf C^{\sslash}$（相应地 $\mathbf C^{\backslash\backslash}$）记作等变化、剪切（相应地反剪切）再去除等变化所得的范畴。最后，若两个范畴都有分次提升 $\mathbf C_{gr},\mathbf D_{gr}$，且存在范畴等价 $\mathbf C_{gr}\simeq\mathbf D_{gr}^{\sslash}$，就称这两个范畴 \emph{Koszul 等价}，记作
\[
\mathbf C\underset{Kos}{\sim}\mathbf D.
\]
\end{defn}

\ysrc{459}

有时用较粗的 Koszul 等价概念会更方便：把上述精细版本过渡到 2-周期范畴即可得到它。

\begin{defn}{A.11}\label{defn:a5-04-5}
设 $\mathbf C$ 是 $k$-线性范畴。$\mathbf C$ 的 \emph{2-周期化}是 2-周期范畴
\[
\mathbf C^{per}:=\mathbf C\otimes_{\Mod(k)}\Mod(k[u,u^{-1}]),
\]
其中 $|u|=2$。
\end{defn}

存在典范函子 $(-)^{per}:\mathbf C\to\mathbf C^{per}$，使得 $\Hom_{\mathbf C^{per}}(X^{per},Y^{per})$ 是链复形 $\Hom_{\mathbf C}(X,Y)$ 的 2-周期化。

由于 Tate 剪切在 2-周期范畴中是恒等，显然 Koszul 等价会在 2-周期范畴上诱导等价
\[
\mathbf C^{per}\simeq\mathbf D^{per}.
\]

\begin{exam}{A.12}\label{exam:a5-04-6}
2-周期化操作会给范畴引入许多新对象。例如，取环面上 ind-凝聚 D-模的范畴 $\breve{\mathcal D}(BT)$；它等价于 $BT$ 的上同调的模范畴 $\Mod(H^\bullet(BT;k))$，后者又等价于 $\Mod(\Sym^\bullet\mathfrak t^*[-2])$。粗略地说，这个范畴在 0 点处只有一个摩天楼对象。然而，其 2-周期化等价于范畴 $\Mod((\Sym^\bullet\mathfrak t^*[-2])[u,u^{-1}])$，后者对每个点 $t\in\mathfrak t$ 都有摩天楼对象，由赋值映射 $u^{-1}\mathrm{ev}_t:\mathfrak t^*[-2]u^{-1}\simeq\mathfrak t^*\to k$ 给出。

我们还注意到，这个范畴的标准分次提升给 $\mathfrak t^*$ 权 $-1$，Tate 剪切把它从次数 2 移到次数 0；相比之下，Rees 参数 $t$ 有权 $1$，Tate 剪切把它变成次数 2 的 $S^1$-形变参数 $u$。
\end{exam}

\subsubsection*{A.1.4.\quad 例}

下面给出 $S^1$-等变范畴的一些非常可算的例子。具体地说，设 $G$ 是阿贝尔仿射代数（可能是形式）群，Cartier 对偶为 $\widehat{G}$。于是经由 Cartier 对偶，有等同
\[
\Coh(\mathcal L(BG))=\Coh(G\times BG)\simeq\Coh(G\times\widehat{G}).
\]
此外，Cartier 对偶给出幺半等价
\[
(\mathrm{QC}(B\Z),\circ)\simeq(\mathrm{QC}(\GG_m),\otimes),
\]
也就是说，范畴上的 $S^1$-作用可以描述为一个 $\mathrm{QC}(\GG_m)$-线性结构。在这一框架中，两者都来自几何，即来自沿典范配对的拉回
\[
G\times\widehat{G}\to\GG_m,\qquad (g,\chi)\mapsto\chi(g).
\]
应用 [Pr11, §3] 的讨论，我们有
\[
\Coh(\mathcal L(BG))^{S^1}\simeq\Coh\bigl((G\times\widehat{G})\times_{\GG_m}\{1\}\bigr),\qquad
\Coh(\mathcal L(BG))^{\mathrm{Tate}}\simeq\mathrm{MF}(G\times\widehat{G},f),
\]
其中 $\mathrm{MF}$ 记矩阵因子化的 2-周期范畴。下面我们在两个具体例子中算出上述内容。

\ysrc{460}

\begin{exam}{A.13}\label{exam:a5-04-7}
取 $X=BT$，其中 $T$ 是代数环面；则 $\mathcal L(BT)\simeq T\times BT$。我们有范畴的等同
\[
\Coh(\mathcal L(BT))\simeq\Coh(T\times BT)\simeq\Coh(T\times X^\bullet(T))
\simeq\bigoplus_{\lambda\in X^\bullet(T)}\Coh(T),
\]
即分解为 $\Rep(T)$ 的同型分量。$S^1$-作用由 $(t,\lambda)\in T\times X^\bullet(T)$ 处的乘以 $t^\lambda$ 给出（见注 A.2），即在 $T\times\{\lambda\}$ 上乘以单项式 $t^\lambda\in\OO(T)$。下面我们计算 $S^1$-等变对象、Tate 局部化与 $S^1$-不变对象（见命题 A.6）。令 $T_\lambda$ 记多项式 $1-t^\lambda$ 的导出零点集。具体地，
\[
T_\lambda=(\OO(T)[\epsilon],d(\epsilon)=1-t^\lambda)\simeq
\begin{cases}
\OO(T)[\epsilon], & \lambda=1,\\
\OO(T)/t^\lambda-1, & \lambda\neq1.
\end{cases}
\]
由注 A.2，对由自同构 $\alpha$ 给出的 $S^1$-作用，其 $S^1$-不变量可以描述为以导出方式附加上方程 $1-\alpha$。于是
\[
\Coh(\mathcal L(BT))^{S^1}\simeq\bigoplus_{\lambda\in X^\bullet(T)}\Coh(T_\lambda).
\]
按 Koszul 对偶 [Pr11, Con. 3.1.5]，我们可以把 $\Coh(\mathcal L(BT))^{S^1}$ 视为 $k[u]$-模范畴，其中 $|u|=2$ 通过上同调算子作用。当 $\lambda\neq1$ 时导出零点集 $T_\lambda$ 光滑，于是任何凝聚层都有有限消解，从而 $u$ 在 $\Coh(T_\lambda)$ 上的作用是挠的，即 $\Coh(T_\lambda)\otimes_{k[u]}k(u)\simeq0$。当 $\lambda=1$ 时，$u$ 自由作用，并且有（另见 [Pr11, Prop. 3.4.1]）：
\[
\Coh(\mathcal L(BT))^{\mathrm{Tate}}\simeq\Coh(T)\otimes\Mod(k[u,u^{-1}]).
\]
特别地，
\[
\Coh(\widehat{\mathcal L}(BT))^{\mathrm{Tate}}
=\Coh(\mathcal L^u(BT))^{\mathrm{Tate}}
\subsetneq\Coh(\mathcal L(BT))^{\mathrm{Tate}},
\]
其中第一个等号成立是因为 $T$ 没有幂幺元素，即 $\mathcal L^u(BT)=\widehat{\mathcal L}(BT)$。而在 $u=0$ 处特殊化，我们得到 [Pr11, Cor. 3.2.4]
\[
\Coh(\mathcal L(BT))^{S^1}\otimes_{k[u]}k
\simeq\bigoplus_{\lambda\in X^\bullet(T)}\Coh_{T_\lambda}(T).
\]
特别地，并非 $\Coh(\mathcal L(BT))$ 的每个对象都是 $S^1$-可等变化的。
\end{exam}

\begin{exam}{A.14}\label{exam:a5-04-8}
我们考虑一个多少正交的例。取 $X=B\GG_a$。于是由 Cartier 对偶有
\[
\Coh(\mathcal L(B\GG_a))\simeq\Coh(\GG_a\times B\GG_a)\simeq\Coh(\GG_a\times\widehat{\GG}_a).
\]
分别在两个 $\GG_a$ 上固定坐标 $x,y$；我们把 $x$ 看作来自 $\OO(\GG_a)$ 的作用，把 $y$ 看作 $1\in\mathrm{Lie}(\GG_a)$ 的作用，即 $V\in\Coh(B\GG_a)$ 上的作用映射是 $t\mapsto\exp(ty)$。于是 $S^1$-作用由恒同函子的自同构 $\exp(xy)$ 给出。取 $S^1$-不变量，就是在导出意义下附加上方程 $\exp(xy)=1$，即 $xy=0$（由于 $y$ 处在 $0$ 的无穷小邻域中，可以忽略 $1$ 的其他对数）。我们有
\[
\Coh(\mathcal L(B\GG_a))^{S^1}\simeq\Mod_{\mathrm{f.g.},\,y\text{-}\mathrm{nil}}(k[x,y]/xy).
\]
Tate 局部化是 $k[x,y]$ 上关于方程 $xy=0$ 的矩阵因子化范畴（其中 $y$ 幂零地作用）。有等价
\[
\Coh(\widehat{\mathcal L}(B\GG_a))^{\mathrm{Tate}}
\simeq\Coh(\mathcal L^u(B\GG_a))^{\mathrm{Tate}}
=\Coh(\mathcal L(B\GG_a))^{\mathrm{Tate}}
\simeq\Mod(k(u)),
\]

\ysrc{461}

其中第一个同构之所以成立，是因为方程 $xy=0$ 的奇点轨迹在原点（因此在 $x=0$ 处完备化不改变矩阵因子化范畴），第二个则来自几何（即 $\mathcal L^u(B\GG_a)=\mathcal L(B\GG_a)$，因为 $\GG_a$ 只含幂幺元素）。最右边的范畴等同是一个计算，例如见 [Dy11, Sec. 5.3]。或者，可以经由 Koszul 对偶（定理 2.6），其中 $\Coh(\widehat{\mathcal L}(B\GG_a))^{\mathrm{Tate}}\simeq\mathcal D(B\GG_a)\otimes\Mod(k(u))\simeq\Mod(k(u))$。$\Coh(\mathcal L(B\GG_a))$ 的所有对象都是 $S^1$-不变的，即
\[
\Coh(\mathcal L(B\GG_a))^{S^1}\otimes_{k[u]}k=\Coh(\mathcal L(B\GG_a)).
\]
由于所有圈都是幂幺的，$S^1$-作用经过一个 $B\GG_a$-作用分解，而后者是 1-仿射的。这个例子容易推广到 $X=BV$ 的情形，其中 $V$ 是交换加法群。
\end{exam}

\subsection*{A.2.\quad Koszul 对偶与等变局部化}

下面我们把注意力转向 Koszul 对偶与等变局部化的讨论中出现的各种技术细节。

\subsubsection*{A.2.1.\quad 重整化范畴}

光滑叠 $X$ 上 D-模的范畴 $\mathcal D(X)$ 通过下降来定义。令 $F\mathcal D(X)$ 记叠 $X$ 上滤过 D-模的范畴，同样通过下降定义。令 $F\mathcal D_c(X)$ 记 $X$ 上凝聚 D-模（带良好滤过）的全子范畴。注意这些一般并不是 $F\mathcal D(X)$ 中的紧对象；我们令 $F\breve{\mathcal D}(X)=\Ind(F\mathcal D_c(X))$ 记 ind-凝聚滤过 D-模的范畴，即关于凝聚 D-模对范畴作重整化。

我们还需要对奇切丛作类似的重整化；讨论见 [CD22] 第 2 节。对任意闭子叠 $Z\subset X$，我们定义 ind-连续凝聚层的范畴 $\breve{\Coh}(\widehat X_Z)$（有时记作 $\breve{\Coh}_Z(X)$）为满足下述条件的对象 $\mathcal F\in\mathrm{QC}^{!}(\widehat X_Z)\simeq\mathrm{QC}^{!}_Z(X)$ 组成的全子范畴：$\mathcal F$ 是 $t$-有界的且几乎 $!$-完美，即对任意集论支撑于 $Z\subset X$ 的闭子叠 $i:Z'\into X$，$!$-限制 $i^{!}\mathcal F$ 具有凝聚上同调。这个范畴具有以下性质（见 [Ch23]）。
\begin{enumerate}[(1)]
\item 它包含沿 $Z\subset X$ 支撑的通常凝聚层范畴，即 $\Coh_Z(X)\subset\breve{\Coh}_Z(X)$。特别地，有 $\Coh(\widehat{\mathcal L}X)=\Coh_{\mathcal L X}(\mathcal L X)\subset\breve{\Coh}(\widehat{\mathcal L}X)$。
\item 令 $\widehat i:\widehat X_Z\into X$ 为形式完备化的包含，则函子 $\widehat i^{!}:\mathrm{QC}^{!}(X)\to\mathrm{QC}^{!}(\widehat X_Z)$ 把 $\Coh(X)$ 送到 $\breve{\Coh}(\widehat X_Z)$。特别地，取 $\widehat z:\widehat{\mathcal L}X\to\mathcal L X$ 为形式圈空间的包含，函子 $\widehat z^{!}:\mathrm{QC}^{!}(\mathcal L X)\to\mathrm{QC}^{!}(\widehat{\mathcal L}X)$ 限制为一个函子 $\Coh(\mathcal L X)\to\breve{\Coh}(\widehat{\mathcal L}X)$。
\item 若 $Z\subset X$ 由幂零理想定义，则 $\Coh(\widehat X_Z)=\breve{\Coh}(\widehat X_Z)=\Coh(X)$。特别地，若 $X$ 是概形，则 $\Coh(\widehat{\mathcal L}X)=\breve{\Coh}(\widehat{\mathcal L}X)=\Coh(\mathcal L X)$。
\item 若 $Z\subset X$ 是导出局部完全交，则 $\mathcal F\in\breve{\Coh}(\widehat X_Z)$ 当且仅当它到 $Z$ 的 $!$-限制属于 $\Coh(Z)$。特别地，若 $p:U\to X$ 是光滑图册，则 $\mathcal F\in\breve{\Coh}(\widehat{\mathcal L}X)$ 当且仅当 $\widehat{\mathcal L}p^{!}\mathcal F\in\Coh(\mathcal L U)$。
\end{enumerate}

我们用 $\widehat{\mathrm{QC}}^{!}(\widehat X_Z):=\Ind(\breve{\Coh}(\widehat X_Z))$ 记其 ind-完备化。为得到到参数处完备化的良好限制函子，这一重整化是必要的。在一般情形下，设 $F:\mathbf C\to\mathbf D$ 是紧生成范畴之间保余极限且保紧对象的函子（即带保余极限右伴随的左伴随），则对紧对象 $X\in\mathbf C$ 有交换图

\ysrc{462}

\[
\begin{tikzcd}[column sep=large,row sep=large]
\End(X)^{\mathrm{op}}\text{-}\mathrm{mod} \arrow[r] \arrow[d, "-\otimes_{\End(X)}\End(F(X))"'] & \mathbf C \arrow[d, "F"]\\
\End(F(X))^{\mathrm{op}}\text{-}\mathrm{mod} \arrow[r] & \mathbf D.
\end{tikzcd}
\]

交换性来自右伴随之间的自动交换性；而 $X$ 的紧性保证 $\Hom(X,-)$ 的左伴随全忠实\footnote{即当 $\Hom(X,-)$ 与余极限交换时，伴随的单位 $M\to\Hom(X,X\otimes_{\End(X)}M)$ 是等价。}（$F(X)$ 同理）。可惜，ind-凝聚层上的 $!$-限制是右伴随，并不保紧对象。另一方面，重整化之后
\[
\begin{aligned}
\widehat\ell^{!}_\alpha\colon\ \mathrm{QC}^{!}(\mathcal L(X/G)) &\longrightarrow\\
&\widehat{\mathrm{QC}}^{!}\bigl(\widehat{\mathcal L}_\alpha(X(\alpha)/G(\alpha))\bigr)
\end{aligned}
\]
成为左伴随，并且按构造保紧对象。特别地，我们有交换方块
\[
\begin{tikzcd}[column sep=large,row sep=large]
\Mod(\mathcal H)\simeq\langle\mathcal S\rangle \arrow[r] \arrow[d, "-\otimes_{\mathcal H}\mathcal H(\alpha)"'] & \mathrm{QC}^{!}(\mathcal L(X/G)) \arrow[d, "\widehat\ell^{!}_\alpha"]\\
\Mod(\mathcal H(\alpha))\simeq\langle\mathcal S(\alpha)\rangle \arrow[r] & \widehat{\mathrm{QC}}^{!}\bigl(\widehat{\mathcal L}_\alpha(X(\alpha)/G(\alpha))\bigr).
\end{tikzcd}
\]

\subsubsection*{A.2.2.\quad 同伦不动点}

定理 2.9 中的等变局部化陈述用到同伦不动点的概念；我们简要概述这一概念，细节请读者参阅 [AKLPR21] 的附录。

\begin{defn}{A.15}\label{defn:a5-04-9}
设 $G$ 是作用于预叠 $X$ 的群预叠。我们用（导出）纤维积定义\emph{同伦不动点} $X^G$：
\[
\begin{tikzcd}[column sep=large,row sep=large]
X^G \arrow[r] \arrow[d] & \Map(BG,X/G) \arrow[d]\\
\{\mathrm{id}_{BG}\} \arrow[r] & \Map(BG,BG).
\end{tikzcd}
\]
\end{defn}

换言之，$X^G$ 是映射 $X/G\to BG$ 的截面预叠。

我们讨论这一构造的若干例子，因为对不同输入的 $G$，它的味道不同。
\begin{enumerate}[(1)]
\item 当 $G$ 线性约化、$X$ 是局部 Noether 的 Artin 叠时，由 [AKLPR21] 的推论 A.35，$X^G$ 的余切复形在 Tor 幅度（即「奇点级别」）上有与 $X$ 的余切复形相同的下界。特别地，若 $X$ 是光滑概形，则由 [AKLPR21] 的命题 A.23，$X^G\subset X$ 是光滑闭子概形，从而 $X^G$ 就是 $X$ 的经典 $G$-不动点。这一观察，结合 $(-)^{hG}$ 与 Zariski 局部化及纤维积交换这一事实，给出一种计算具有给定表示的导出概形的同伦不动点的方法。
\item 若 $G$ 非约化，$X^G$ 可能有导出结构。例如取 $G=\GG_a$ 平凡地作用于 $X$，则 $X^G=\mathcal L^uX$ 是幂幺圈空间；当 $X$ 是概形时，它就是通常的圈空间 $\mathcal L X$。
\end{enumerate}

\ysrc{463}

\begin{enumerate}[(1),start=3]
\item 假定 $G$ 约化，但不假定 $X$ 光滑，此时 $X^G$ 仍可能有非平凡的导出结构。例如，幂零锥 $\mathcal N=\mathfrak g\times_{\mathfrak g\sslash G}\{0\}$ 是非导出的完全交。然而取 $G$-不变量后得到 $\mathcal N^G=\mathfrak z(\mathfrak g)\times_{\mathfrak g\sslash G}\{0\}$，只要 $\mathfrak g$ 不交换（即 $G$ 不是环面），它就有导出结构。当 $G$ 是环面、$X$ 拟光滑且经典时，可以论证 $X^G$ 是经典的。\footnote{可以仿照定理 2.9 的论证，归约到 $X$ 是光滑 $G$-概形的纤维积的情形，再归约到光滑 $G$-概形导出交的情形，然后借助切空间论证。}
\item 若 $G$ 是拓扑群、平凡地作用于 $X$，则同伦不动点 $X^G=\Map(BG,X)=\mathcal Loc_G(X)$ 是 $X$ 上 $G$-局部系统的导出模叠。这个叠常常有非平凡的导出结构，即使 $X$ 本身没有。
\item 若 $G=\Z$ 经由自同构 $\phi:X\to X$ 作用于 $X$，则得到 $\phi$ 在 $X$ 上的导出不动点 $\mathcal L_\phi X$。特别地，即使 $X$ 是光滑概形，只要不动点的维数不为 $0$，它也是导出的。
\item 更一般地，若 $G$ 是离散循环群（从而是交换群），生成元为 $\eta\in G$，则存在典范等价 $X^G\simeq\mathcal L_\eta(X/G)$。例如可以这样看出：在下图中外侧方块与左侧方块都是笛卡儿的：
\[
\begin{tikzcd}[column sep=large,row sep=large]
X^G \arrow[r] \arrow[d] & \Map(BG,X/G) \arrow[r] \arrow[d] & \mathcal L(X/G) \arrow[d]\\
\{e\} \arrow[r] & \Hom(BG,BG) \arrow[r] & \mathcal L(BG),
\end{tikzcd}
\]
并且注意底部的复合恰是映射 $\{\eta\}\to G/G$。
\item 设 $G=\Z/n\Z$ 平凡地作用于概形 $X$。则 $X^G$ 是映射 $\mathbb T_X[-1]\to\mathbb T_X[-1]$ 在 $X$ 上的相对核，该映射由 $S^1$ 上的 $n$ 次映射（即乘以 $n$）诱导，当 $n$ 在 $k$ 中可逆时它是同构。注意 $G=\Z/n\Z$ 在 $k$ 上线性约化，当且仅当 $n$ 在 $k$ 中可逆。
\end{enumerate}

\subsubsection*{A.2.3.\quad 中心平移与扭变}

定理 2.9 中的局部化映射把对 $\mathcal L(X/G)$ 在 $[z]\in G\sslash G$ 上的研究，归结为对 $\mathcal L(X(\alpha)/G(\alpha))$ 在 $[\alpha]\in G(\alpha)\sslash G(\alpha)$ 上的研究。后者的好处在于：$\alpha\in G(\alpha)$ 是中心元，并且平凡地作用于 $X(\alpha)$；因此我们可以作一次中心平移，把 $\alpha$ 平移为单位元。这次平移对圈旋转并非 $S^1$-等变的。\footnote{例如，$\mathcal L(BG)=G/G$ 上的圈旋转在 $g\in G/G$ 处的作用由共轭 $g$ 的自同构给出，而中心平移并不保持它。}本节的目标是描述这次平移，以及使之为等变的一个扭 $S^1$-作用。

我们先描述这次平移。定理 2.9 中构造的局部化映射分解为两步：
\[
\mathcal L(X(\alpha)/G(\alpha))\to\mathcal L(X/G(\alpha))\to\mathcal L(X/G),
\]
其中第一个映射给出局部化，第二个是 $\alpha\in G\sslash G$ 附近的简单 étale 基变换。我们将假定已经作过这次基变换，只专注于第一个映射，从而可以假定 $\alpha\in Z(G)$ 是中心元。我们需要引入 $\alpha$-平凡化的概念。

\ysrc{464}

\begin{defn}{A.16}\label{defn:a5-04-10}
设 $G$ 是线性代数群，其单位分支约化，作用于预叠 $X$；设 $H\subset G$ 是正规子群，记 $G'=G/H$。$G$ 在 $X$ 上作用的\emph{$H$-平凡化}，指 $X$ 上的一个 $G'$-作用以及一个等价
\[
X/G\simeq X/G'\times_{BG'}BG.
\]
当 $H=\langle\alpha\rangle$ 且 $\alpha\in Z(G)$ 为中心元时，我们把 $H$-平凡化称为\emph{$\alpha$-平凡化}。
\end{defn}

\begin{prop}{A.17}\label{prop:a5-04-11}
设 $G$ 作用于导出概形 $X$，$A\subset G$ 是中心子群。同伦不动点 $X^A$ 带有标准的 $A$-平凡化。
\end{prop}

\begin{proof}
首先注意：若 $i:BA\to BG$ 记 $A$ 到 $G$ 的包含所定义的映射，则有典范等价
\[
X^A\simeq\{i\}\times_{\Map(BA,BG)}\Map(BA,X/G).
\]
在右端，由于 $A$ 是中心子群且两者都约化，包含 $\{i\}\into\Map(BA,BG)\simeq\Hom_{grp}(A,G)/G$ 可提升为一个映射 $\{i\}/G\into\Map(BG,BG)$。由于 $A$ 平凡地作用于 $\Map(BA,BG)$，这个映射还下降为映射 $\{i\}/G'\to\Hom_{grp}(A,G)/G$。这些数据给 $X^A$ 提供了所要的结构。
\end{proof}

再次令 $\alpha\in G$ 为中心元，取 $A=\langle\alpha\rangle$，并设 $X$ 是带 $A$-平凡化的导出概形。下面我们定义 $\mathcal L(X/G)$ 上沿 $\alpha$ 的平移映射
\[
sh_\alpha:\mathcal L(X/G)\to\mathcal L(X/G),
\]
它与乘以 $\alpha$ 的映射 $\mu_\alpha:\mathcal L(BG)=G/G\to\mathcal L(BG)=G/G$ 相容，定义如下。典范的 $\alpha$-平凡化 $X^\alpha/G\simeq X^\alpha/G'\times_{BG'}BG$ 诱导典范等价 $\mathcal L(X^\alpha/G)\simeq\mathcal L(X^\alpha/G')\times_{G'/G'}G/G$，而自同构由右端的「$G/G$ 上乘以 $z$」映射与别处的恒同给出。在定理 2.9 的框架中，我们得到一个复合，称为（幂幺）$\alpha$-移位局部化映射
\[
s\ell^u_\alpha:\ \mathcal L^u(X(\alpha)/G(\alpha))
\xrightarrow[sh_\alpha]{\sim}\mathcal L^u_\alpha(X(\alpha)/G(\alpha))
\xrightarrow[\ell^u_\alpha]{\sim}\mathcal L^u_\alpha(X/G).
\]

左边的平移映射对圈旋转并非 $S^1$-等变；下面引入 $S^1$-作用的扭变。

对 $\alpha\in Z(G)$，乘以 $\alpha$ 定义一个群同态 $\Z\times G\to G$，从而给出作用映射 $S^1\times BG\to BG$。若 $G$ 作用于一个 $\alpha$-平凡化的概形 $X$，则这个 $S^1$-作用给出 $X/G\simeq X/G'\times_{BG'}BG$ 上的 $S^1$-作用，其中取 $G'=G/\langle\alpha\rangle$，并令 $S^1$ 平凡地作用于 $X/G'$ 与 $BG'$。我们称所得的 $S^1$-作用为 $X/G$ 上的\emph{$\alpha$-扭 $S^1$-作用} $\sigma_\alpha$。由于 $S^1$-作用 $\mathcal L(\sigma_\alpha)$ 与圈旋转 $\rho$ 交换，我们定义\emph{$\alpha$-移位圈旋转}，记作 $\rho(\alpha)$ 或 $S^1(\alpha)$，为 $S^1\times S^1$-作用 $\rho\times\mathcal L(\sigma_\alpha)$ 的对角作用。移位局部化映射关于这一扭变等变。

\begin{cor}{A.18}\label{cor:a5-04-12}
移位局部化映射给出等价
\[
s\ell^u_\alpha:\mathcal L^u(X(\alpha)/G(\alpha))\isom\mathcal L^u_\alpha(X/G),
\]
它在源上关于 $\rho(\alpha)$、在靶上关于 $\rho$ 是 $S^1$-等变的；对完备化与特殊化的局部化映射的平移 $s\ell^{\wedge}_\alpha$ 与 $s\ell'_\alpha$ 同样如此。
\end{cor}

\ysrc{465}

\subsubsection*{A.2.4.\quad 扭作用的平凡块}

为应用第 2.3 节讨论的 Koszul 对偶（它涉及未扭作用的 $S^1$-等变性），我们希望找出一类层构成的子范畴：这些层位于半单参数 $\alpha$ 处的导出圈空间上，并使 $\alpha$-扭变成为平凡。这很有用，因为幂幺圈空间上的 $\rho$ 圆作用经过一个 $B\GG_a$-作用分解，而扭后的 $\rho(\alpha)$ 作用则不然（因为它有非平凡的半单部分）。这个障碍使得我们无法用 Koszul 对偶把 $\Coh(\widehat{\mathcal L}_\alpha(X/G))^{S^1}$ 与某类 D-模范畴等同起来，而第 2.4.2 节正是这样做的。

为定义这个子范畴，我们给上面讨论的几何 $\alpha$-扭 $S^1$-作用 $\sigma(\alpha)$ 一个范畴论解释。

\begin{defn}{A.19}\label{defn:a5-04-13}
设 $G$ 是仿射代数群，$A\subset Z(G)$ 是中心子群，$G'=G/A$，$\mathbf C$ 是 $\mathrm{QC}(BG)$-模范畴。$\mathbf C$ 的\emph{$A$-平凡化}，指如下数据：一个 $\mathrm{QC}(BG')$-模范畴 $\mathbf C'$，以及一个等同 $\mathbf C\simeq\mathbf C'\otimes_{\mathrm{QC}(BG')}\mathrm{QC}(BG)$。我们定义\emph{$A$-平凡对象}的子范畴 $\mathbf C_A\subset\mathbf C$ 为自然函子 $\mathbf C'\to\mathbf C$ 的本质像。当 $A=\langle\alpha\rangle$ 且 $\alpha\in Z(G)$ 时，记作 $\mathbf C_\alpha$。
\end{defn}

进一步，若 $G$ 约化，则 $\alpha$-平凡对象构成 $\mathbf C$ 的一个直和项（而不只是一个子范畴）。由于 $G$ 约化，其中心半单；令 $X^\bullet(A)$ 记 $A$ 的特征标群，则有分解 $\mathrm{QC}(BG)=\Rep(G)=\bigoplus_{\chi\in X^\bullet(A)}\Rep(G)_\chi$，其中平凡同型对应于 $\Rep(G')$。于是 $\mathbf C$ 的 $\alpha$-平凡化给出直和分解
\[
\mathbf C\simeq\bigoplus_{\chi\in X^\bullet(A)}\mathbf C_\chi
\simeq\bigoplus_{\chi\in X^\bullet(A)}\mathbf C'\otimes_{\Rep(G')}\Rep(G)_\chi.
\]
我们将考察自然界中出现的下列 $\alpha$-平凡化范畴。
\begin{enumerate}[(1)]
\item 若 $X/G$ 是整体商叠，则叠的 $\alpha$-平凡化给出范畴 $\mathrm{QC}(X)$、$\mathrm{QC}^{!}(X)$ 等的 $z$-平凡化\tn{底本此处作「$z$-trivialization」，依上下文（第 465 页）应为「$\alpha$-trivialization」。}。
\item 若 $\alpha\in G$ 是中心元，则 $\mathcal L(BG)=G/G$ 有典范的 $\alpha$-平凡化。
\item 结合上面两条，$\mathcal L(X/G)$ 上有 $\alpha$-平凡化。
\end{enumerate}

\begin{exam}{A.20}\label{exam:a5-04-14}
回到例 A.13。回忆 $T_\lambda=\ker(\lambda)$ 是 $\lambda\in X^\bullet(T)$ 的（导出）核，其中 $\alpha\in T$。我们有描述
\[
\Coh(\mathcal L_\alpha(BT))^{S^1}=\bigoplus_{\lambda\in X^\bullet(T)}\Coh(T_\lambda),\qquad
\Coh(\mathcal L(BT))^{S^1(\alpha)}=\bigoplus_{\lambda\in X^\bullet(T)}\Coh(\alpha^{-1}T_\lambda),
\]
而 $\alpha$-平凡块恰对应于那些满足 $\alpha\in T_\lambda$ 的项。特别地，应用 Tate 构造，$\alpha$-平凡块的包含诱导等价
\[
\Coh(\mathcal L(BT))^{\mathrm{Tate}}\simeq\Coh(\mathcal L(BT))^{\mathrm{Tate}(\alpha)}
\]
对所有扭 $\alpha$ 成立。
\end{exam}

\ysrc{466}

\subsubsection*{A.2.5.\quad 非约化群的局部化}

定理 2.9 中的等变局部化可以推广到非约化群 $K$，办法如下：给定商叠 $X/K$，我们可以转而考虑商叠 $(X\times^K G)/G$，并在 $G\sslash G$ 上作等变局部化。一般而言，映射 $K\sslash K\to G\sslash G$ 既不是单射也不是满射（例如 $K=U$ 是幂幺子群时）。

我们特别关注 $K=B\subset G$ 是 $G$ 的 Borel 子群的情形。此时有图
\[
\begin{tikzcd}[column sep=large,row sep=large]
\mathcal L(X/B)=\mathcal L\bigl((X\times^B G)/G\bigr) \arrow[r] & \mathcal L(BB)=B/B \arrow[r] \arrow[d] & \mathcal L(BG)=G/G \arrow[d]\\
& B\sslash B=H \arrow[r,"\nu"] & G\sslash G=H\sslash W.
\end{tikzcd}
\]

对给定的 $\alpha\in H$ 及 $[\alpha]=\nu(\alpha)\in H\sslash W$，我们可以把 $G/B$ 的 $\alpha$-不动点与 $\nu^{-1}([\alpha])\times G(\alpha)/B(\alpha)$ 作 $W_{G(\alpha)}$-等变的等同，其中 $G(\alpha)$ 是 $\alpha$ 的中心化子，$B(\alpha)\subset G(\alpha)$ 是 Borel 子群。于是
\[
\mathcal L_{[\alpha]}(BB)=\mathcal L_{[\alpha]}(G\backslash G/B)
=\mathcal L_{[\alpha]}\bigl(G(\alpha)\backslash(G/B)^\alpha\bigr)
=\coprod_{\nu^{-1}([\alpha])}\mathcal L_{[\alpha]}\bigl(G(\alpha)\backslash G(\alpha)/B(\alpha)\bigr).
\]
我们定义 $\mathcal L_\alpha(BB)$ 为对应于 $\alpha\in\nu^{-1}([\alpha])$ 的连通分支，然后对 $\alpha\in B\sslash B=H$ 可以定义
\[
\mathcal L_\alpha(X/B):=\mathcal L_{[\alpha]}\bigl((X\times^B G)/G\bigr)
\times_{\mathcal L_{[\alpha]}(BB)}\mathcal L_\alpha(BB).
\]
细节见 [Ch20] 的例 1.0.6。

\section*{参考文献}

\begin{refs}
\item[{[ABPS17]}] Anne-Marie Aubert, Paul Baum, Roger Plymen, and Maarten Solleveld, \emph{Conjectures about $p$-adic groups and their noncommutative geometry}, Around Langlands correspondences, Contemp. Math., vol. 691, Amer. Math. Soc., Providence, RI, 2017, pp. 15--51. MR3666049
\item[{[ABV92]}] Jeffrey Adams, Dan Barbasch, and David A. Vogan Jr., \emph{The Langlands classification and irreducible characters for real reductive groups}, Progress in Mathematics, vol. 104, Birkh\"auser Boston, Inc., Boston, MA, 1992, DOI 10.1007/978-1-4612-0383-4. MR1162533
\item[{[AKLPR21]}] Dhyan Aranha, Adeel A. Khan, Alexei Latyntsev, Hyeonjun Park, and Charanya Ravi, \emph{Virtual localization revisited. part A}, Adv. Math. \textbf{479} (2025), no. part A, Paper No. 110434, DOI 10.1016/j.aim.2025.110434. MR4933128
\item[{[Ber16]}] Joseph Bernstein, \emph{Stacks in representation theory}, Tsinghua lectures in mathematics, Adv. Lect. Math. (ALM), vol. 45, Int. Press, Somerville, MA, 2019, pp. 123--139. MR3971550
\item[{[Bez06]}] Roman Bezrukavnikov, \emph{Noncommutative counterparts of the Springer resolution}, International Congress of Mathematicians. Vol. II, Eur. Math. Soc., Z\"urich, 2006, pp. 1119--1144. MR2275638
\item[{[Bez16]}] Roman Bezrukavnikov, \emph{On two geometric realizations of an affine Hecke algebra}, Publ. Math. Inst. Hautes \'Etudes Sci. \textbf{123} (2016), 1--67, DOI 10.1007/s10240-015-0077-x. MR3502096
\item[{[BGS96]}] Alexander Beilinson, Victor Ginzburg, and Wolfgang Soergel, \emph{Koszul duality patterns in representation theory}, J. Amer. Math. Soc. \textbf{9} (1996), no. 2, 473--527, DOI 10.1090/S0894-0347-96-00192-0. MR1322847
\item[{[BCHN22]}] David Ben-Zvi, Harrison Chen, David Helm, and David Nadler, \emph{Coherent Springer theory and the categorical Deligne-Langlands correspondence}, Invent. Math. \textbf{235} (2024), no. 2, 255--344, DOI 10.1007/s00222-023-01224-2. MR4689369
\item[{[BCHN]}] David Ben-Zvi, Harrison Chen, David Helm, and David Nadler, \emph{Equivariant localization for Langlands parameters}, in preparation.

\ysrc{467}

\item[{[BFN10]}] David Ben-Zvi, John Francis, and David Nadler, \emph{Integral transforms and Drinfeld centers in derived algebraic geometry}, J. Amer. Math. Soc. \textbf{23} (2010), no. 4, 909--966, DOI 10.1090/S0894-0347-10-00669-7. MR2669705
\item[{[BN07]}] David Ben-Zvi and David Nadler, \emph{Loop spaces and Langlands parameters}, Preprint, arXiv:0706.0322.
\item[{[BN10a]}] David Ben-Zvi and David Nadler, \emph{Loop spaces and connections}, J. Topol. \textbf{5} (2012), no. 2, 377--430, DOI 10.1112/jtopol/jts007. MR2928082
\item[{[BN10b]}] David Ben-Zvi and David Nadler, \emph{Loop spaces and representations}, Duke Math. J. \textbf{162} (2013), no. 9, 1587--1619, DOI 10.1215/00127094-2266130. MR3079256
\item[{[BN12]}] David Ben-Zvi and David Nadler, \emph{Beilinson-Bernstein localization over the Harish-Chandra center}, arXiv:1209.0188.
\item[{[BN16]}] David Ben-Zvi and David Nadler, \emph{Betti geometric Langlands}, Algebraic geometry: Salt Lake City 2015, Proc. Sympos. Pure Math., vol. 97, Amer. Math. Soc., Providence, RI, 2018, pp. 3--41, DOI 10.1090/pspum/097.2/01. MR3821166
\item[{[BNP17]}] David Ben-Zvi, David Nadler, and Anatoly Preygel, \emph{A spectral incarnation of affine character sheaves}, Compos. Math. \textbf{153} (2017), no. 9, 1908--1944, DOI 10.1112/S0010437X17007278. MR3705281
\item[{[BV21]}] Roman Bezrukavnikov and Kari Vilonen, \emph{Koszul duality for quasi-split real groups}, Invent. Math. \textbf{226} (2021), no. 1, 139--193, DOI 10.1007/s00222-021-01050-4. MR4309492
\item[{[BY13]}] Roman Bezrukavnikov and Zhiwei Yun, \emph{On Koszul duality for Kac-Moody groups}, Represent. Theory \textbf{17} (2013), 1--98, DOI 10.1090/S1088-4165-2013-00421-1. MR3003920
\item[{[CD22]}] Harrison Chen and Gurbir Dhillon, \emph{A Langlands dual realization of coherent sheaves on the nilpotent cone}, arXiv:2310.10539v1, 2023.
\item[{[CG97]}] Neil Chriss and Victor Ginzburg, \emph{Representation theory and complex geometry}, Birkh\"auser Boston, Inc., Boston, MA, 1997. MR1433132
\item[{[Ch20]}] Harrison Chen, \emph{Equivariant localization and completion in cyclic homology and derived loop spaces}, Adv. Math. \textbf{364} (2020), 107005, 56, DOI 10.1016/j.aim.2020.107005. MR4060042
\item[{[Ch23]}] Harrison Chen, \emph{Categorical cyclic homology and filtered D-modules on stacks: Koszul duality}, arXiv:2301.06949v1, 2023.
\item[{[DHKM20]}] Jean-Fran\c{c}ois Dat, David Helm, Robert Kurinczuk, and Gilbert Moss, \emph{Moduli of Langlands parameters}, J. Eur. Math. Soc. (JEMS) \textbf{27} (2025), no. 5, 1827--1927, DOI 10.4171/jems/1599. MR4889237
\item[{[Dy11]}] Tobias Dyckerhoff, \emph{Compact generators in categories of matrix factorizations}, Duke Math. J. \textbf{159} (2011), no. 2, 223--274, DOI 10.1215/00127094-1415869. MR2824483
\item[{[EH14]}] Matthew Emerton and David Helm, \emph{The local Langlands correspondence for GLn in families} (English, with English and French summaries), Ann. Sci. \'Ec. Norm. Sup\'er. (4) \textbf{47} (2014), no. 4, 655--722, DOI 10.24033/asens.2224. MR3250061
\item[{[Fa16]}] Laurent Fargues, \emph{Geometrization of the local Langlands correspondence: an overview}, arXiv:1602.00999, 2016.
\item[{[FS21]}] Laurent Fargues and Peter Scholze, \emph{Geometrization of the local Langlands correspondence}, arXiv:2102.13459, 2021.
\item[{[Ga14]}] Dennis Gaitsgory, \emph{Sheaves of categories and the notion of 1-affineness}, Stacks and categories in geometry, topology, and algebra, Contemp. Math., vol. 643, Amer. Math. Soc., Providence, RI, 2015, pp. 127--225, DOI 10.1090/conm/643/12899. MR3381473
\item[{[Ga16]}] Dennis Gaitsgory, \emph{From geometric to function-theoretic Langlands (or how to invent shtukas)}, arXiv:1606.09608v1, 2016.
\item[{[GKRV20]}] D. Gaitsgory, D. Kazhdan, N. Rozenblyum, and Y. Varshavsky, \emph{A toy model for the Drinfeld-Lafforgue shtuka construction}, Indag. Math. (N.S.) \textbf{33} (2022), no. 1, 39--189, DOI 10.1016/j.indag.2021.11.002. MR4368480
\item[{[GR17]}] Dennis Gaitsgory and Nick Rozenblyum, \emph{A study in derived algebraic geometry. Vol. I. Correspondences and duality}, Mathematical Surveys and Monographs, vol. 221, American Mathematical Society, Providence, RI, 2017, DOI 10.1090/surv/221.1. MR3701352
\item[{[GL18]}] Alain Genestier and Vincent Lafforgue, \emph{Chtoucas restreints pour les groupes r\'eductifs et param\'etrisation de Langlands locale}, arXiv:1709.00978v2, 2018.

\ysrc{468}

\item[{[GKM98]}] Mark Goresky, Robert Kottwitz, and Robert MacPherson, \emph{Equivariant cohomology, Koszul duality, and the localization theorem}, Invent. Math. \textbf{131} (1998), no. 1, 25--83, DOI 10.1007/s002220050197. MR1489894
\item[{[H20]}] David Helm, \emph{Curtis homomorphisms and the integral Bernstein center for GLn}, Algebra Number Theory \textbf{14} (2020), no. 10, 2607--2645, DOI 10.2140/ant.2020.14.2607. MR4190413
\item[{[He20]}] Eugen Hellmann, \emph{On the derived category of the Iwahori-Hecke algebra}, Compos. Math. \textbf{159} (2023), no. 5, 1042--1110, DOI 10.1112/s0010437x23007145. MR4586564
\item[{[HM18]}] David Helm and Gilbert Moss, \emph{Converse theorems and the local Langlands correspondence in families}, Invent. Math. \textbf{214} (2018), no. 2, 999--1022, DOI 10.1007/s00222-018-0816-y. MR3867634
\item[{[HL22]}] Quoc P. Ho and Penghui Li, \emph{Revisiting mixed geometry}, Preprint, arXiv:2202.04833.
\item[{[Ka91]}] M. M. Kapranov, \emph{On DG-modules over the de Rham complex and the vanishing cycles functor}, Algebraic geometry (Chicago, IL, 1989), Lecture Notes in Math., vol. 1479, Springer, Berlin, 1991, pp. 57--86, DOI 10.1007/BFb0086264. MR1181207
\item[{[K89]}] Masaki Kashiwara, \emph{Representation theory and D-modules on flag varieties: Orbits unipotentes et repr\'esentations, III}, Ast\'erisque \textbf{173-174} (1989), 9, 55--109. MR1021510
\item[{[K08]}] Masaki Kashiwara, \emph{Equivariant derived category and representation of real semisimple Lie groups}, Representation theory and complex analysis, Lecture Notes in Math., vol. 1931, Springer, Berlin, 2008, pp. 137--234, DOI 10.1007/978-3-540-76892-0\_3. MR2409699
\item[{[KL87]}] David Kazhdan and George Lusztig, \emph{Proof of the Deligne-Langlands conjecture for Hecke algebras}, Invent. Math. \textbf{87} (1987), no. 1, 153--215, DOI 10.1007/BF01389157. MR862716
\item[{[KS94]}] Masaki Kashiwara and Wilfried Schmid, \emph{Quasi-equivariant D-modules, equivariant derived category, and representations of reductive Lie groups}, Lie theory and geometry, Progr. Math., vol. 123, Birkh\"auser Boston, Boston, MA, 1994, pp. 457--488, DOI 10.1007/978-1-4612-0261-5\_16. MR1327544
\item[{[La18a]}] Vincent Lafforgue, \emph{Chtoucas pour les groupes r\'eductifs et param\'etrisation de Langlands globale} (French), J. Amer. Math. Soc. \textbf{31} (2018), no. 3, 719--891, DOI 10.1090/jams/897. MR3787407
\item[{[La18b]}] Vincent Lafforgue, \emph{Shtukas for reductive groups and Langlands correspondence for function fields}, Proceedings of the International Congress of Mathematicians---Rio de Janeiro 2018. Vol. I. Plenary lectures, World Sci. Publ., Hackensack, NJ, 2018, pp. 635--668. MR3966741
\item[{[LaZ19]}] Vincent Lafforgue and Xinwen Zhu, \emph{D\'ecomposition au-dessus des param\`etres de Langlands elliptiques}, arXiv:1811.07976v2, 2019.
\item[{[Lur1]}] Jacob Lurie, \emph{Higher topos theory}, Annals of Mathematics Studies, vol. 170, Princeton University Press, Princeton, NJ, 2009, DOI 10.1515/9781400830558. MR2522659
\item[{[Lur2]}] Jacob Lurie, \emph{Higher algebra}, http://math.ias.edu/\textasciitilde lurie/papers/HA.pdf.
\item[{[Lur3]}] Jacob Lurie, \emph{Rotation invariance in algebraic K-theory}, available at http://math.ias.edu/\textasciitilde lurie/ (2015).
\item[{[Lu88]}] George Lusztig, \emph{Cuspidal local systems and graded Hecke algebras. I}, Inst. Hautes \'Etudes Sci. Publ. Math. \textbf{67} (1988), 145--202. MR972345
\item[{[Lu89]}] George Lusztig, \emph{Affine Hecke algebras and their graded version}, J. Amer. Math. Soc. \textbf{2} (1989), no. 3, 599--635, DOI 10.2307/1990945. MR991016
\item[{[Lu95a]}] George Lusztig, \emph{Cuspidal local systems and graded Hecke algebras. II}, with errata for Part I [Inst. Hautes \'Etudes Sci. Publ. Math. No. 67 (1988), 145--202; MR0972345 (90e:22029)], Representations of groups (Banff, AB, 1994), CMS Conf. Proc., vol. 16, Amer. Math. Soc., Providence, RI, 1995, pp. 217--275. MR1357201
\item[{[Lu95b]}] George Lusztig, \emph{Classification of unipotent representations of simple $p$-adic groups}, Internat. Math. Res. Notices \textbf{11} (1995), 517--589, DOI 10.1155/S1073792895000353. MR1369407

\ysrc{469}

\item[{[LuY20]}] George Lusztig and Zhiwei Yun, \emph{Endoscopy for Hecke categories, character sheaves and representations}, Forum Math. Pi \textbf{8} (2020), e12, 93, DOI 10.1017/fmp.2020.9. MR4108915
\item[{[MB25]}] L. Mason-Brown, \emph{Arthur's conjectures and the orbit method for real reductive groups}, The Langlands program: I. Representations, Shimura varieties, and shtukas, Proc. Sympos. Pure Math., vol. 112.1, Amer. Math. Soc., Providence, RI, 2025.
\item[{[MUV92]}] I. Mirkovi\'c, T. Uzawa, and K. Vilonen, \emph{Matsuki correspondence for sheaves}, Invent. Math. \textbf{109} (1992), no. 2, 231--245, DOI 10.1007/BF01232026. MR1172690
\item[{[Pr11]}] Anatoly Preygel, \emph{Thom-Sebastiani and duality for matrix factorizations}, arXiv:1101.5834v1, 2011.
\item[{[Pr15]}] Anatoly Preygel, \emph{Ind-coherent complexes on loop spaces and connections}, Stacks and categories in geometry, topology, and algebra, Contemp. Math., vol. 643, Amer. Math. Soc., Providence, RI, 2015, pp. 289--323, DOI 10.1090/conm/643/12902. MR3381476
\item[{[Ri13]}] Laura Rider, \emph{Formality for the nilpotent cone and a derived Springer correspondence}, Adv. Math. \textbf{235} (2013), 208--236, DOI 10.1016/j.aim.2012.12.001. MR3010057
\item[{[RiRu16]}] Laura Rider and Amber Russell, \emph{Perverse sheaves on the nilpotent cone and Lusztig's generalized Springer correspondence}, Lie algebras, Lie superalgebras, vertex algebras and related topics, Proc. Sympos. Pure Math., vol. 92, Amer. Math. Soc., Providence, RI, 2016, pp. 273--292, DOI 10.1090/pspum/092/01592. MR3644235
\item[{[RiRu21]}] Laura Rider and Amber Russell, \emph{Formality and Lusztig's generalized Springer correspondence}, Algebr. Represent. Theory \textbf{24} (2021), no. 3, 699--714, DOI 10.1007/s10468-020-09966-w. MR4651747
\item[{[SV17]}] Yiannis Sakellaridis and Akshay Venkatesh, \emph{Periods and harmonic analysis on spherical varieties} (English, with English and French summaries), Ast\'erisque \textbf{396} (2017), viii+360. MR3764130
\item[{[So01]}] Wolfgang Soergel, \emph{Langlands' philosophy and Koszul duality}, Algebra---representation theory (Constanta, 2000), NATO Sci. Ser. II Math. Phys. Chem., vol. 28, Kluwer Acad. Publ., Dordrecht, 2001, pp. 379--414. MR1858045
\item[{[Sol18]}] Maarten Solleveld, \emph{A local Langlands correspondence for unipotent representations}, Amer. J. Math. \textbf{145} (2023), no. 3, 673--719. MR4596175
\item[{[Sol19]}] Maarten Solleveld, \emph{On unipotent representations of ramified $p$-adic groups}, Represent. Theory \textbf{27} (2023), 669--716, DOI 10.1090/ert/652. MR4620884
\item[{[Sol22a]}] Maarten Solleveld, \emph{Graded Hecke algebras, constructible sheaves and the $p$-adic Kazhdan-Lusztig conjecture}, J. Algebra \textbf{667} (2025), 865--910, DOI 10.1016/j.jalgebra.2025.01.006. MR4854590
\item[{[Sol22b]}] Maarten Solleveld, \emph{Graded Hecke algebras and equivariant constructible sheaves on the nilpotent cone}, Q. J. Math. \textbf{76} (2025), no. 1, 109--146, DOI 10.1093/qmath/haae065. MR4865035
\item[{[TV]}] Bertrand To\"en and Gabriele Vezzosi, \emph{Alg\`ebres simpliciales $S^1$-\'equivariantes, th\'eorie de de Rham et th\'eor\`emes HKR multiplicatifs} (French, with English and French summaries), Compos. Math. \textbf{147} (2011), no. 6, 1979--2000, DOI 10.1112/S0010437X11005501. MR2862069
\item[{[Vo82]}] David A. Vogan Jr., \emph{Irreducible characters of semisimple Lie groups. IV. Character-multiplicity duality}, Duke Math. J. \textbf{49} (1982), no. 4, 943--1073. MR683010
\item[{[Vo93]}] David A. Vogan Jr., \emph{The local Langlands conjecture}, Representation theory of groups and algebras, Contemp. Math., vol. 145, Amer. Math. Soc., Providence, RI, 1993, pp. 305--379, DOI 10.1090/conm/145/1216197. MR1216197
\item[{[Zh15]}] Xinwen Zhu, \emph{The geometric Satake correspondence for ramified groups} (English, with English and French summaries), Ann. Sci. \'Ec. Norm. Sup\'er. (4) \textbf{48} (2015), no. 2, 409--451, DOI 10.24033/asens.2248. MR3346175
\item[{[Zh17]}] Xinwen Zhu, \emph{An introduction to affine Grassmannians and the geometric Satake equivalence}, Geometry of moduli spaces and representation theory, IAS/Park City Math. Ser., vol. 24, Amer. Math. Soc., Providence, RI, 2017, pp. 59--154. MR3752460
\item[{[Zh18]}] Xinwen Zhu, \emph{Geometric Satake, categorical traces, and arithmetic of Shimura varieties}, Current developments in mathematics 2016, Int. Press, Somerville, MA, 2018, pp. 145--206. MR3837875

\ysrc{470}

\item[{[Zhu25a]}] Xinwen Zhu, \emph{Coherent sheaves on the stack of Langlands parameters}, The Langlands program: II. Geometrization of the Langlands correspondence, Proc. Sympos. Pure Math., vol. 112.2, Amer. Math. Soc., Providence, RI, 2025.
\item[{[Zhu25b]}] Xinwen Zhu, \emph{Tame categorical local Langlands correspondence}, arXiv:2504.07482, 2025.\tn{底本在第 470 页把两篇不同文献都标作 [Zhu25]；今仿第 4 篇参考文献的处理，分别记作 [Zhu25a]、[Zhu25b]。}
\end{refs}

\medskip
\begin{center}
{\small Department of Mathematics, University of Texas, Austin, Texas 78712-0257\\[0.3em]
\itshape Email address: \texttt{benzvi@math.utexas.edu}}

\medskip
{\small Institute of Mathematics, Academia Sinica, Taipei 106319, Taiwan\\[0.3em]
\itshape Email address: \texttt{chenhi.math@gmail.com}}

\medskip
{\small Department of Mathematics, Imperial College, London SW7 2BU, United Kingdom\\[0.3em]
\itshape Email address: \texttt{dhelm@imperial.ac.uk}}

\medskip
{\small Department of Mathematics, University of California, Berkeley, California 94720-3840\\[0.3em]
\itshape Email address: \texttt{nadler@math.berkeley.edu}}
\end{center}

% 本片止于印刷页 470（PDF 483），第 5 篇全文结束（参考文献自印 466 [ABPS17] 起，至印 470 [Zhu25b]；
% 其后为四位作者的通讯地址）。下一篇为第 6 篇（Feng–Harris, Derived structures in the Langlands
% correspondence），其篇首页即 PDF 484 = 印刷页 471，故本片不涉及该篇任何内容。
